\documentclass[11pt,a4paper]{article}
\usepackage[utf8]{inputenc}
\usepackage[T1]{fontenc}
\usepackage{amsmath,amssymb,amsthm}
\usepackage{listings}
\usepackage{geometry}
\usepackage{hyperref}
\geometry{margin=2.5cm}

\newtheorem{lemma}{Lemma}

\lstset{
  language=C++,
  basicstyle=\ttfamily\small,
  keywordstyle=\bfseries,
  breaklines=true,
  numbers=left,
  numberstyle=\tiny,
  frame=single,
  tabsize=2
}

\title{IOI 2019 -- Split the Attractions}
\author{Solution}
\date{}

\begin{document}
\maketitle

\section{Problem Summary}
Given a connected graph with $n$ nodes and $m$ edges, partition the nodes into three groups of sizes $a$, $b$, $c$ ($a + b + c = n$) such that each group induces a connected subgraph. Output the assignment or report impossibility.

\section{Solution}

\subsection{Centroid-Based Approach}
\begin{enumerate}
  \item Sort sizes so $a \le b \le c$. Build a spanning tree and find its centroid.
  \item At the centroid, every subtree has size $\le n/2$. Since $a \le n/3 \le n/2$, we can carve out a connected subgraph of size exactly $a$.
  \item If some child subtree has size $\ge a$, recurse into it to extract exactly $a$ nodes via DFS.
  \item Otherwise, greedily merge entire subtrees until the total reaches $a$.
  \item After removing the $a$-group, the remaining $n - a$ nodes form a connected tree (containing the centroid). Repeat to extract a group of size $b$; the remainder forms the $c$-group.
\end{enumerate}

\begin{lemma}
A spanning tree with centroid $v$ always allows extracting a connected subtree of any size $s \le n/2$.
\end{lemma}
\begin{proof}
Each child subtree of $v$ has size $\le n/2$. If any subtree has size $\ge s$, we can greedily take $s$ nodes from it by DFS. Otherwise, merge subtrees: since each is $< s$ and their total is $n - 1 \ge s$, we can accumulate exactly $s$ by adding whole subtrees (the overshoot is at most one subtree size $< s$, so we recurse into the last subtree).
\end{proof}

\section{Implementation}

\begin{lstlisting}
#include <bits/stdc++.h>
using namespace std;

vector<int> find_split(int n, int a, int b, int c,
                       vector<int> p, vector<int> q) {
    int m = p.size();
    vector<vector<int>> adj(n);
    for (int i = 0; i < m; i++) {
        adj[p[i]].push_back(q[i]);
        adj[q[i]].push_back(p[i]);
    }

    // Sort sizes; track original labels
    int sizes[3] = {a, b, c}, labels[3] = {1, 2, 3};
    for (int i = 0; i < 3; i++)
        for (int j = i + 1; j < 3; j++)
            if (sizes[i] > sizes[j]) {
                swap(sizes[i], sizes[j]);
                swap(labels[i], labels[j]);
            }

    // BFS spanning tree from 0
    vector<int> par(n, -1);
    vector<bool> vis(n, false);
    vector<vector<int>> tree(n);
    vector<int> order;
    queue<int> bfs;
    bfs.push(0); vis[0] = true;
    while (!bfs.empty()) {
        int u = bfs.front(); bfs.pop(); order.push_back(u);
        for (int v : adj[u]) if (!vis[v]) {
            vis[v] = true; par[v] = u;
            tree[u].push_back(v);
            bfs.push(v);
        }
    }

    // Compute subtree sizes and find centroid
    vector<int> sz(n, 1);
    for (int i = (int)order.size() - 1; i >= 0; i--) {
        int u = order[i];
        for (int v : tree[u]) sz[u] += sz[v];
    }
    int centroid = 0;
    for (int u : order)
        for (int v : tree[u])
            if (sz[v] > n / 2) { centroid = v; }

    // Re-root at centroid
    fill(vis.begin(), vis.end(), false);
    for (int i = 0; i < n; i++) tree[i].clear();
    order.clear();
    bfs.push(centroid); vis[centroid] = true;
    while (!bfs.empty()) {
        int u = bfs.front(); bfs.pop(); order.push_back(u);
        for (int v : adj[u]) if (!vis[v]) {
            vis[v] = true; tree[u].push_back(v); bfs.push(v);
        }
    }
    fill(sz.begin(), sz.end(), 1);
    for (int i = (int)order.size() - 1; i >= 0; i--) {
        int u = order[i];
        for (int v : tree[u]) sz[u] += sz[v];
    }

    vector<int> ans(n, -1);

    // Extract exactly 'need' nodes from subtree of u, label them
    function<int(int, int)> extract = [&](int u, int need) -> int {
        if (need <= 0) return 0;
        ans[u] = 0;  // temporary mark
        int taken = 1;
        for (int v : tree[u]) {
            if (taken >= need) break;
            int take = min(sz[v], need - taken);
            taken += extract(v, take);
        }
        return taken;
    };

    int small = sizes[0];
    extract(centroid, small);  // temporarily mark with 0

    // Assign label to marked nodes
    int cnt1 = 0;
    for (int i = 0; i < n; i++)
        if (ans[i] == 0) { ans[i] = labels[0]; cnt1++; }

    // Extract second group from remaining connected component
    // Find a remaining node connected to centroid (or centroid itself if unmarked)
    int start2 = -1;
    if (ans[centroid] != labels[0]) {
        start2 = centroid;
    } else {
        for (int v : tree[centroid])
            if (ans[v] != labels[0]) { start2 = v; break; }
    }

    if (start2 != -1) {
        // BFS on remaining nodes to extract sizes[1]
        vector<bool> vis2(n, false);
        queue<int> bfs2;
        bfs2.push(start2); vis2[start2] = true;
        int taken = 0;
        while (!bfs2.empty() && taken < sizes[1]) {
            int u = bfs2.front(); bfs2.pop();
            if (ans[u] != -1) continue;
            ans[u] = labels[1]; taken++;
            for (int v : adj[u])
                if (!vis2[v] && ans[v] == -1) {
                    vis2[v] = true; bfs2.push(v);
                }
        }
    }

    // Remaining nodes get third label
    for (int i = 0; i < n; i++)
        if (ans[i] == -1) ans[i] = labels[2];

    return ans;
}
\end{lstlisting}

\section{Complexity Analysis}
\begin{itemize}
  \item \textbf{Time}: $O(n + m)$ -- spanning tree, centroid, and extractions are all linear.
  \item \textbf{Space}: $O(n + m)$.
\end{itemize}

\end{document}
